% ============================================================
% UNIDAD 5: EQUIPOS DE REEMPLAZO (RETROFIT)
% ============================================================
\clearpage
\section*{UNIDAD 5: EQUIPOS DE REEMPLAZO}
\subsection*{G5, GI 275, G3X Touch, G500/G600 TXi, AXIS y Aspen}

\begin{recuadro}
\textbf{Objetivo de la unidad:} conocer los instrumentos digitales que se instalan para
modernizar aeronaves existentes, sus componentes y sensores, su forma de instalación y sus
diferencias, para poder recomendarlos, instalarlos y mantenerlos.
\end{recuadro}

% ------------------------------------------------------------
\section*{1. POR QUÉ MODERNIZAR UN PANEL}

La mayoría de los aviones de la aviación general argentina tienen décadas de antigüedad e
instrumentos de vacío. Reemplazarlos por instrumentos digitales: elimina la bomba de vacío (una de
las fallas más comunes y peligrosas en vuelo por instrumentos); mejora la confiabilidad; agrega
funciones (HSI con GPS, alertas, motor digital, ADS-B); y aumenta el valor de la aeronave. Los
equipos de reemplazo están pensados para \textbf{instalarse en los orificios existentes} del panel
o con modificaciones mínimas, mediante un STC que cubre muchos modelos.

% ------------------------------------------------------------
\section*{2. INSTRUMENTOS AUTÓNOMOS}

\subsubsection*{2.1. Garmin G5}
\begin{itemize}
  \item Pantalla de 3,5 pulgadas que se monta en un \textbf{orificio estándar de 3 1/8 pulgadas}.
  \item Contiene \textbf{AHRS y ADC internos}: se conecta a las líneas pitot y estática.
  \item Puede configurarse como \textbf{indicador de actitud} (con velocidad y altitud en cintas) o
        como \textbf{indicador de rumbo / HSI} (con el magnetómetro remoto GMU 11 y, para mostrar
        la navegación de un GPS o NAV existente, un adaptador GAD 29).
  \item \textbf{Batería de respaldo} interna: sigue funcionando si se pierde la alimentación.
  \item En aeronaves certificadas se instala por STC como reemplazo del horizonte artificial y del
        indicador de rumbo; también se usa como \textbf{instrumento de reserva} en paneles más
        grandes. Hay una versión experimental para aeronaves amateur y LSA.
  \item Se integra con el piloto automático GFC 500.
\end{itemize}

\subsubsection*{2.2. Garmin GI 275}
\begin{itemize}
  \item También se instala en un orificio de 3 1/8 pulgadas, con pantalla táctil y perilla.
  \item Según la versión y la configuración reemplaza: \textbf{horizonte (ADI)}, \textbf{HSI},
        \textbf{CDI}, \textbf{indicadores de motor (EIS)} o funciona como \textbf{pequeño MFD}
        (mapa, tránsito, meteorología).
  \item La versión ADI tiene AHRS y ADC internos; puede tener batería de respaldo.
  \item Permite reemplazar instrumento por instrumento manteniendo el formato del panel, y es
        compatible con muchos pilotos automáticos antiguos.
\end{itemize}

\subsubsection*{2.3. Aspen Evolution (EFD 1000 PFD)}
\begin{itemize}
  \item Pantalla vertical de 6 pulgadas que ocupa el lugar del \textbf{horizonte y del indicador de
        rumbo}, montada sobre los dos orificios existentes: se instala con muy poca modificación
        del panel.
  \item PFD completo (actitud, cintas de velocidad y altitud, HSI) con sensores internos y un
        \textbf{módulo de sensores remoto (RSM)} montado en el fuselaje (magnetómetro y temperatura).
  \item Batería interna de respaldo. Se conecta a navegadores GPS y radios existentes.
  \item Se complementa con una o dos pantallas MFD de la misma familia.
\end{itemize}

% ------------------------------------------------------------
\section*{3. SISTEMAS DE PANTALLAS}

\subsubsection*{3.1. Garmin G3X Touch}
\begin{itemize}
  \item Pantallas táctiles de 7 y 10,6 pulgadas, configurables como PFD, MFD o ambos.
  \item Sensores remotos: \textbf{ADAHRS} (GSU 25), magnetómetro (GMU 11), sonda de temperatura,
        unidad de motor (GEA 24); se comunican por un \textbf{bus CAN}.
  \item Nació para aeronaves experimentales y LSA; luego se aprobó por STC para muchos modelos
        certificados. Se combina con el G5 como reserva y con el piloto automático GFC 500.
\end{itemize}

\subsubsection*{3.2. Garmin G500 / G600 TXi}
\begin{itemize}
  \item Pantallas táctiles (10,6 y 7 pulgadas) para aeronaves certificadas, en formato horizontal
        o vertical. PFD, MFD e indicación de motor opcional.
  \item Trabajan junto a los navegadores GPS de la familia GTN y radios externas, que siguen en la
        columna central del panel.
\end{itemize}

\subsubsection*{3.3. Garmin AXIS}
\begin{itemize}
  \item Familia de pantallas presentada por Garmin en julio de 2026 para aviones con motor a pistón
        (monomotores y bimotores livianos), mediante un STC con lista de modelos, y para aeronaves
        experimentales y LSA.
  \item Su característica principal es que la pantalla \textbf{integra el navegador GPS IFR, las
        radios de comunicaciones y navegación y el panel de audio}: se elimina la columna central de
        equipos, como en un G1000, pero como modernización.
  \item Comparte muchos sensores, unidades remotas, recortes del panel y agujeros de montaje con el
        G3X Touch, lo que facilita la actualización desde ese sistema.
  \item Agrega conectividad inalámbrica (Bluetooth y Wi-Fi) para bases de datos y planes de vuelo.
\end{itemize}

% ------------------------------------------------------------
\section*{4. CUADRO COMPARATIVO}
\begin{center}
\footnotesize
\begin{tabular}{|>{\raggedright\arraybackslash}p{0.13\linewidth}|>{\raggedright\arraybackslash}p{0.19\linewidth}|>{\raggedright\arraybackslash}p{0.21\linewidth}|>{\raggedright\arraybackslash}p{0.15\linewidth}|>{\raggedright\arraybackslash}p{0.16\linewidth}|}
\hline
\textbf{Equipo} & \textbf{Tipo} & \textbf{Instalación} & \textbf{Sensores} & \textbf{Uso típico} \\
\hline
G5 & Instrumento autónomo & Orificio de 3 1/8 pulg. & Internos + GMU 11 & Reemplazo de vacío;
reserva \\ \hline
GI 275 & Instrumento autónomo multifunción & Orificio de 3 1/8 pulg. & Internos (ADI) & Reemplazo
gradual, motor \\ \hline
Aspen EFD 1000 & PFD autónomo & Dos orificios (horizonte y rumbo) & Internos + RSM & PFD con
poca obra \\ \hline
G3X Touch & Sistema de pantallas & Recorte del panel & ADAHRS remoto, CAN & Experimental, LSA y
certificados \\ \hline
G500/G600 TXi & Sistema de pantallas & Recorte del panel & ADAHRS remoto & Certificados con
columna de radios \\ \hline
AXIS & Pantallas con radios integradas & Recorte del panel & Remotos (familia G3X) & Panel
completo nuevo \\ \hline
\end{tabular}
\end{center}
Otras marcas presentes en el mercado: Dynon (SkyView, con versiones certificadas por STC), uAvionix
(AV-30), Avidyne, MGL y otras. El criterio de elección incluye: aeronave y su STC, uso (VFR o IFR),
piloto automático, presupuesto, peso, consumo eléctrico y soporte técnico disponible en el país.

% ------------------------------------------------------------
\section*{5. ACTIVIDADES DE LA UNIDAD}
\begin{enumerate}
  \item Proponer dos alternativas de modernización para un Cessna 172 de 1978 con instrumentos de
        vacío, una económica y una completa, con un esquema del panel resultante.
  \item Leer el manual de instalación del G5 y listar: conexiones eléctricas, conexiones de pitot y
        estática, requisitos del magnetómetro y calibraciones necesarias.
  \item Comparar el diagrama de interconexión de un G3X Touch con el de un G1000: ¿qué buses usa
        cada uno?
  \item Investigar qué equipos de esta unidad tienen STC aceptado por la ANAC y cómo se tramita
        esa aceptación.
\end{enumerate}
